\documentclass[12pt,a4paper]{article}

\usepackage[brazilian]{babel}

\usepackage{amsmath}
\usepackage{amssymb}
\usepackage{enumitem}
\usepackage{etoolbox}
\usepackage{tocloft}
\usepackage{tikz}
\usepackage[most]{tcolorbox}
\usepackage{xcolor}

\setlength{\cftsecnumwidth}{3.2cm}
\setlength{\cftsubsecnumwidth}{1.5cm}

\renewcommand{\thesection}{Capítulo \arabic{section} -}
\renewcommand{\thesubsection}{\arabic{section}.\arabic{subsection}}
\pretocmd{\subsection}{\clearpage}{}{}

% Cria o comando \qedsymbolcolorido com um quadrado de 1.8ex x 1.8ex
% Você pode ajustar a cor (ex: blue!70!black, teal, black, etc.)
\newcommand{\fimresolucao}{\hfill\textcolor{blue!70!black}{\rule{1.8ex}{1.8ex}}}

% Caixa para ferramentas e recursos teóricos
\newtcolorbox{recursosbox}[1][]{
    enhanced,
    colback=blue!3!white,
    colframe=blue!60!black,
    coltitle=white,
    fonttitle=\bfseries,
    title=Ferramentas e Recursos da Seção,
    arc=2mm,
    boxrule=0.8pt,
    left=6pt, right=6pt, top=6pt, bottom=6pt,
    #1
}

% Caixa para alertas de erros frequentes e dicas práticas
\newtcolorbox{dicasbox}[1][]{
    enhanced,
    colback=orange!4!white,
    colframe=orange!80!black,
    coltitle=white,
    fonttitle=\bfseries,
    title=Atenção aos Detalhes e Dicas de Resolução,
    arc=2mm,
    boxrule=0.8pt,
    left=6pt, right=6pt, top=6pt, bottom=6pt,
    #1
}

\title{
    \Huge \textbf{Um Curso de Cálculo Volume 1} \\
    \vspace{0.3cm}
    \Large 6\textordfeminine{} Edição (Guidorizzi) \\
    \vspace{0.5cm}
    \large \textit{Resoluções Completas e Passo a Passo do Capítulo 8 Derivadas}
}
\author{Rogério Campos}

\begin{document}

\maketitle
\newpage

\tableofcontents
\newpage

\setcounter{section}{6}
\setcounter{subsection}{0}

\section{Derivadas}

\subsection{Introdução}
Esta subseção não possui exercícios

\subsection{Derivada de uma função}

% ==========================================================
% FERRAMENTAS E RECURSOS TEÓRICOS (GUIDORIZZI)
% ==========================================================
\begin{recursosbox}
\textbf{Definição Formal de Derivada num Ponto:}
Sejam $f$ uma função e $p$ um ponto de seu domínio. O limite
\[
\lim_{x \to p} \frac{f(x) - f(p)}{x - p}
\]
quando existe e é finito, denomina-se \textbf{derivada de $f$ em $p$} e indica-se por $f'(p)$. Se $f$ admite derivada em $p$, dizemos que $f$ é \textbf{derivável} (ou diferenciável) em $p$.

\vspace{0.2cm}
\textbf{Formas Equivalentes de Cálculo:}
Fazendo a mudança de variável $x = p + h$ (onde $x \to p \iff h \to 0$), obtém-se a forma alternativa clássica por acréscimo:
\begin{itemize}[noitemsep, topsep=2pt]
    \item \textbf{Forma por aproximação pontual:}
    \[
    f'(p) = \lim_{x \to p} \frac{f(x) - f(p)}{x - p}
    \]
    \item \textbf{Forma por acréscimo ($h = \Delta x$):}
    \[
    f'(p) = \lim_{h \to 0} \frac{f(p + h) - f(p)}{h}
    \]
\end{itemize}

\vspace{0.2cm}
\textbf{Interpretação Geométrica e Reta Tangente:}
\begin{itemize}[noitemsep, topsep=2pt]
    \item $f'(p)$ representa geometricamente o \textbf{coeficiente angular} ($m$) da reta tangente ao gráfico de $f$ no ponto $(p, f(p))$.
    \item A equação da reta tangente ao gráfico de $f$ em $(p, f(p))$ é dada por:
    \[
    y - f(p) = f'(p)(x - p)
    \]
\end{itemize}
\end{recursosbox}

% ==========================================================
% ATENÇÃO AOS DETALHES E DICAS PRÁTICAS
% ==========================================================
\begin{dicasbox}
\textbf{Dificuldades Frequentes e Estratégias de Resolução:}
\begin{itemize}[noitemsep, topsep=2pt]
    \item \textbf{Escolha da forma do limite:} Para pontos numéricos fixos (ex.: $p = 2$), a forma $\frac{f(x) - f(p)}{x - p}$ frequentemente simplifica fatorações polinomiais mais rápido. Já para encontrar a função derivada geral $f'(x)$, a forma por acréscimo com $h \to 0$ costuma ser mais direta.
    \item \textbf{Cuidado com o quociente indeterminante:} O limite da derivada sempre produz a forma $\left[\frac{0}{0}\right]$. Você \textbf{deve} manipular algebricamente a expressão (fatorando, racionalizando ou cancelando termos) até eliminar o fator $(x - p)$ ou o denominador $h$ antes de avaliar o limite.
    \item \textbf{Funções por várias sentenças ou módulos:} Para verificar derivabilidade em pontos de transição (ex.: em $x = 0$ para $f(x) = |x|$), é obrigatório calcular os \textbf{limites laterais} do quociente de Newton:
    \[
    f'_+(p) = \lim_{x \to p^+} \frac{f(x) - f(p)}{x - p} \quad \text{e} \quad f'_-(p) = \lim_{x \to p^-} \frac{f(x) - f(p)}{x - p}
    \]
    $f'(p)$ existirá se, e somente se, $f'_+(p) = f'_-(p)$ e ambos forem finitos.
    \item \textbf{Derivabilidade implica Continuidade:} Se $f$ é derivável em $p$, então $f$ é automaticamente contínua em $p$. Se a função apresentar descontinuidade em $p$, a derivada em $p$ não existirá imediatamente (sem necessidade de calcular o limite).
\end{itemize}
\end{dicasbox}
\vspace{0.5cm}


\vspace{0.5cm}
\noindent\textbf{Aquecendo os motores}

\vspace{0.5cm}
\noindent\textbf{Exercício 1.} Seja $f(x) = x^2 + 1$. Calcule:

\begin{enumerate}
    \renewcommand{\theenumi}{\alph{enumi}}
    \renewcommand{\labelenumi}{\theenumi)}
    \item $f'(x)$
    \begin{recursosbox}[title=Recursos Utilizados]
    Utilizaremos a definição de derivada por acréscimo:
    \[
    f'(x) = \lim_{h \to 0} \frac{f(x + h) - f(x)}{h}
    \]
    \textbf{Ferramentas Algébricas:}
    \begin{itemize}[noitemsep, topsep=2pt]
        \item Expansão de binômio: $(x + h)^2 = x^2 + 2hx + h^2$.
        \item Fatoração do quociente incremental para simplificar o fator comum $h$.
    \end{itemize}
    \end{recursosbox}

    Sabemos pela definição formal que:
    \begin{equation}
        f'(x) = \lim_{h \to 0} \frac{f(x + h) - f(x)}{h} \label{eq:def_derivada}
    \end{equation}

    Primeiro, manipulamos o quociente incremental:
    \begin{equation}
        \frac{f(x + h) - f(x)}{h} \label{eq:quociente}
    \end{equation}

    Como foi dado $f(x) = x^2 + 1$, temos:
    \[
    f(x + h) = (x + h)^2 + 1 = x^2 + 2hx + h^2 + 1
    \]
    Substituindo $f(x + h)$ e $f(x)$ em \eqref{eq:quociente} para $h \neq 0$:
    \begin{align*}
        \frac{f(x + h) - f(x)}{h} &= \frac{(x^2 + 2hx + h^2 + 1) - (x^2 + 1)}{h} \\[0.25cm]
        &= \frac{2hx + h^2}{h} \\[0.25cm]
        &= \frac{h(2x + h)}{h} \\[0.25cm]
        &= 2x + h
    \end{align*}

    Daí, aplicando o limite em \eqref{eq:def_derivada} quando $h \to 0$:
    \[
    f'(x) = \lim_{h \to 0} (2x + h) = 2x
    \]

    Portanto:
    \[
    f'(x) = 2x 
    \]
    \fimresolucao

    \item $f'(0)$
    \setcounter{equation}{0}

    \begin{recursosbox}[title=Recursos Utilizados]
    Utilizaremos a definição formal de derivada no ponto $p = 0$:
    \[
    f'(0) = \lim_{x \to 0} \frac{f(x) - f(0)}{x - 0}
    \]
    \end{recursosbox}

    Pela definição formal de derivada no ponto $p = 0$, temos:
    \begin{equation}
        f'(0) = \lim_{x \to 0} \frac{f(x) - f(0)}{x - 0} \label{eq:def_derivada_ponto0}
    \end{equation}

    Primeiro, manipulamos o quociente incremental para $x \neq 0$:
    \begin{equation}
        \frac{f(x) - f(0)}{x - 0} \label{eq:quociente_ponto0}
    \end{equation}

    Sendo $f(x) = x^2 + 1$, calculamos o valor da função em $p = 0$:
    \[
    f(0) = 0^2 + 1 = 1
    \]

    Substituindo $f(x)$ e $f(0)$ em \eqref{eq:quociente_ponto0}:
    \begin{align*}
        \frac{f(x) - f(0)}{x - 0} &= \frac{(x^2 + 1) - 1}{x} \\[0.25cm]
        &= \frac{x^2}{x} \\[0.25cm]
        &= x
    \end{align*}

    Daí, aplicando o limite em \eqref{eq:def_derivada_ponto0} quando $x \to 0$:
    \[
    f'(0) = \lim_{x \to 0} x = 0
    \]

    Portanto:
    \[
    f'(0) = 0
    \]
    \fimresolucao

    \item $f'(1)$
    \setcounter{equation}{0}

    \begin{recursosbox}[title=Recursos Utilizados]
    Utilizaremos a definição formal de derivada no ponto $p = 1$:
    \[
    f'(1) = \lim_{x \to 1} \frac{f(x) - f(1)}{x - 1}
    \]
    \textbf{Ferramentas Algébricas:}
    \begin{itemize}[noitemsep, topsep=2pt]
        \item Fatoração por diferença de dois quadrados: $x^2 - 1 = (x - 1)(x + 1)$.
    \end{itemize}
    \end{recursosbox}

    Pela definição formal de derivada no ponto $p = 1$, temos:
    \begin{equation}
        f'(1) = \lim_{x \to 1} \frac{f(x) - f(1)}{x - 1} \label{eq:def_derivada_ponto1}
    \end{equation}

    Primeiro, manipulamos o quociente incremental para $x \neq 1$:
    \begin{equation}
        \frac{f(x) - f(1)}{x - 1} \label{eq:quociente_ponto1}
    \end{equation}

    Sendo $f(x) = x^2 + 1$, calculamos o valor da função em $p = 1$:
    \[
    f(1) = 1^2 + 1 = 2
    \]

    Substituindo $f(x)$ e $f(1)$ em \eqref{eq:quociente_ponto1}:
    \begin{align*}
        \frac{f(x) - f(1)}{x - 1} &= \frac{(x^2 + 1) - 2}{x - 1} \\[0.25cm]
        &= \frac{x^2 - 1}{x - 1} \\[0.25cm]
        &= \frac{(x - 1)(x + 1)}{x - 1} \\[0.25cm]
        &= x + 1
    \end{align*}

    Daí, aplicando o limite em \eqref{eq:def_derivada_ponto1} quando $x \to 1$:
    \[
    f'(1) = \lim_{x \to 1} (x + 1) = 1 + 1 = 2
    \]

    Portanto:
    \[
    f'(1) = 2
    \]
    \fimresolucao
\end{enumerate}

\vspace{0.5cm}
\noindent\textbf{Exercício 2.}

\vspace{0.5cm}
\noindent\textbf{Exercício 3.}
\begin{enumerate}
    \renewcommand{\theenumi}{\alph{enumi}}
    \renewcommand{\labelenumi}{\theenumi)}
    \item 
    \item 
    \item 
\end{enumerate}

\vspace{0.5cm}
\noindent\textbf{Exercício 4.}
\begin{enumerate}
    \renewcommand{\theenumi}{\alph{enumi}}
    \renewcommand{\labelenumi}{\theenumi)}
    \item 
    \item 
    \item 
    \item 
    \item 
    \item 
    \item 
    \item 
\end{enumerate}

\vspace{0.5cm}
\noindent\textbf{Exercício 5.}
\begin{enumerate}
    \renewcommand{\theenumi}{\alph{enumi}}
    \renewcommand{\labelenumi}{\theenumi)}
    \item 
    \item 
    \item 
    \item 
\end{enumerate}

\vspace{0.5cm}
\noindent\textbf{Exercício 6.}
\begin{enumerate}
    \renewcommand{\theenumi}{\alph{enumi}}
    \renewcommand{\labelenumi}{\theenumi)}
    \item 
    \item 
    \item 
    \item 
    \item 
    \item 
    \item 
    \item 
\end{enumerate}

\vspace{0.5cm}
\noindent\textbf{Exercício 7.}

\vspace{0.5cm}
\noindent\textbf{Exercício 8.}

\vspace{0.5cm}
\noindent\textbf{Exercício 9.}

\vspace{0.5cm}
\noindent\textbf{Exercício 10.}

\vspace{0.5cm}
\noindent\textbf{Exercício 11.}

\vspace{0.5cm}
\noindent\textbf{Exercício 12.}

\vspace{0.5cm}
\noindent\textbf{Exercício 13.}

\vspace{0.5cm}
\noindent\textbf{Exercício 14.}

\vspace{0.5cm}
\noindent\textbf{Exercício 15.}
\begin{enumerate}
    \renewcommand{\theenumi}{\alph{enumi}}
    \renewcommand{\labelenumi}{\theenumi)}
    \item 
    \item 
\end{enumerate}

\vspace{0.5cm}
\noindent\textbf{Exercício 16.}
\begin{enumerate}
    \renewcommand{\theenumi}{\alph{enumi}}
    \renewcommand{\labelenumi}{\theenumi)}
    \item 
    \item 
\end{enumerate}

\vspace{0.5cm}
\noindent\textbf{Exercício 17.}
\begin{enumerate}
    \renewcommand{\theenumi}{\alph{enumi}}
    \renewcommand{\labelenumi}{\theenumi)}
    \item 
    \item 
\end{enumerate}

\vspace{0.5cm}
\noindent\textbf{Exercício 18.}

\vspace{0.5cm}
\noindent\textbf{Exercício 19.}

\subsection{Derivadas de $x^n$ e $\sqrt[n]{x}$}

\vspace{0.5cm}
\noindent\textbf{Exercício 1.}
\begin{enumerate}
    \renewcommand{\theenumi}{\alph{enumi}}
    \renewcommand{\labelenumi}{\theenumi)}
    \item 
    \item 
    \item 
\end{enumerate}

\vspace{0.5cm}
\noindent\textbf{Exercício 2.}
\begin{enumerate}
    \renewcommand{\theenumi}{\alph{enumi}}
    \renewcommand{\labelenumi}{\theenumi)}
    \item 
    \item 
    \item 
    \item 
    \item 
    \item 
    \item 
    \item 
\end{enumerate}

\vspace{0.5cm}
\noindent\textbf{Exercício 3.}

\vspace{0.5cm}
\noindent\textbf{Exercício 4.}

\vspace{0.5cm}
\noindent\textbf{Exercício 5.}
\begin{enumerate}
    \renewcommand{\theenumi}{\alph{enumi}}
    \renewcommand{\labelenumi}{\theenumi)}
    \item 
    \item 
    \item 
\end{enumerate}

\vspace{0.5cm}
\noindent\textbf{Exercício 6.}
\begin{enumerate}
    \renewcommand{\theenumi}{\alph{enumi}}
    \renewcommand{\labelenumi}{\theenumi)}
    \item 
    \item 
    \item 
    \item 
\end{enumerate}

\vspace{0.5cm}
\noindent\textbf{Exercício 7.}

\vspace{0.5cm}
\noindent\textbf{Exercício 8.}

\vspace{0.5cm}
\noindent\textbf{Exercício 9.}

\subsection{Derivadas de $e^x$ e $\ln(x)$}

\vspace{0.5cm}
\noindent\textbf{Exercício 1.}

\vspace{0.5cm}
\noindent\textbf{Exercício 2.}

\vspace{0.5cm}
\noindent\textbf{Exercício 3.}

\vspace{0.5cm}
\noindent\textbf{Exercício 4.}
\begin{enumerate}
    \renewcommand{\theenumi}{\alph{enumi}}
    \renewcommand{\labelenumi}{\theenumi)}
    \item 
    \item 
    \item 
    \item 
\end{enumerate}

\vspace{0.5cm}
\noindent\textbf{Exercício 5.}

\vspace{0.5cm}
\noindent\textbf{Exercício 6.}
\begin{enumerate}
    \renewcommand{\theenumi}{\alph{enumi}}
    \renewcommand{\labelenumi}{\theenumi)}
    \item 
    \item 
    \item 
    \item 
\end{enumerate}

\subsection{Derivadas de Funções Trigonométricas}

\vspace{0.5cm}
\noindent\textbf{Exercício 1.}
\begin{enumerate}
    \renewcommand{\theenumi}{\alph{enumi}}
    \renewcommand{\labelenumi}{\theenumi)}
    \item 
    \item 
\end{enumerate}

\vspace{0.5cm}
\noindent\textbf{Exercício 2.}

\vspace{0.5cm}
\noindent\textbf{Exercício 3.}
\begin{enumerate}
    \renewcommand{\theenumi}{\alph{enumi}}
    \renewcommand{\labelenumi}{\theenumi)}
    \item 
    \item 
    \item 
    \item 
\end{enumerate}

\vspace{0.5cm}
\noindent\textbf{Exercício 4.}
\begin{enumerate}
    \renewcommand{\theenumi}{\alph{enumi}}
    \renewcommand{\labelenumi}{\theenumi)}
    \item 
    \item 
\end{enumerate}

\vspace{0.5cm}
\noindent\textbf{Exercício 5.}

\vspace{0.5cm}
\noindent\textbf{Exercício 6.}
\begin{enumerate}
    \renewcommand{\theenumi}{\alph{enumi}}
    \renewcommand{\labelenumi}{\theenumi)}
    \item 
    \item 
\end{enumerate}

\vspace{0.5cm}
\noindent\textbf{Exercício 7.}
\begin{enumerate}
    \renewcommand{\theenumi}{\alph{enumi}}
    \renewcommand{\labelenumi}{\theenumi)}
    \item 
    \item 
\end{enumerate}

\subsection{Derivabilidade e Continuidade}

\vspace{0.5cm}
\noindent\textbf{Exercício 1.}
\begin{enumerate}
    \renewcommand{\theenumi}{\alph{enumi}}
    \renewcommand{\labelenumi}{\theenumi)}
    \item 
    \item 
\end{enumerate}

\vspace{0.5cm}
\noindent\textbf{Exercício 2.}
\begin{enumerate}
    \renewcommand{\theenumi}{\alph{enumi}}
    \renewcommand{\labelenumi}{\theenumi)}
    \item 
    \item 
\end{enumerate}

\vspace{0.5cm}
\noindent\textbf{Exercício 3.}
\begin{enumerate}
    \renewcommand{\theenumi}{\alph{enumi}}
    \renewcommand{\labelenumi}{\theenumi)}
    \item 
    \item 
\end{enumerate}

\subsection{Regras de Derivação}

\vspace{0.5cm}
\noindent\textbf{Exercício 1.}
\begin{enumerate}
    \renewcommand{\theenumi}{\alph{enumi}}
    \renewcommand{\labelenumi}{\theenumi)}
    \item 
    \item 
    \item 
    \item 
    \item 
    \item 
    \item 
    \item 
    \item 
    \item 
    \item 
    \item 
    \item 
\end{enumerate}

\vspace{0.5cm}
\noindent\textbf{Exercício 2.}

\vspace{0.5cm}
\noindent\textbf{Exercício 3.}
\begin{enumerate}
    \renewcommand{\theenumi}{\alph{enumi}}
    \renewcommand{\labelenumi}{\theenumi)}
    \item 
    \item 
\end{enumerate}

\vspace{0.5cm}
\noindent\textbf{Exercício 4.}
\begin{enumerate}
    \renewcommand{\theenumi}{\alph{enumi}}
    \renewcommand{\labelenumi}{\theenumi)}
    \item 
    \item 
    \item 
\end{enumerate}

\vspace{0.5cm}
\noindent\textbf{Exercício 5.}

\vspace{0.5cm}
\noindent\textbf{Exercício 6.}
\begin{enumerate}
    \renewcommand{\theenumi}{\alph{enumi}}
    \renewcommand{\labelenumi}{\theenumi)}
    \item 
    \item 
    \item 
\end{enumerate}

\vspace{0.5cm}
\noindent\textbf{Exercício 7.}
\begin{enumerate}
    \renewcommand{\theenumi}{\alph{enumi}}
    \renewcommand{\labelenumi}{\theenumi)}
    \item 
    \item 
    \item 
    \item 
    \item 
    \item 
    \item 
    \item 
\end{enumerate}

\vspace{0.5cm}
\noindent\textbf{Exercício 8.}
\begin{enumerate}
    \renewcommand{\theenumi}{\alph{enumi}}
    \renewcommand{\labelenumi}{\theenumi)}
    \item 
    \item 
    \item 
    \item 
\end{enumerate}

\vspace{0.5cm}
\noindent\textbf{Exercício 9.}
\begin{enumerate}
    \renewcommand{\theenumi}{\alph{enumi}}
    \renewcommand{\labelenumi}{\theenumi)}
    \item 
    \item 
    \item 
    \item 
    \item 
    \item 
    \item 
    \item 
    \item 
    \item 
    \item 
    \item 
    \item 
    \item 
    \item 
    \item 
    \item 
    \item 
\end{enumerate}

\vspace{0.5cm}
\noindent\textbf{Exercício 10.}
\begin{enumerate}
    \renewcommand{\theenumi}{\alph{enumi}}
    \renewcommand{\labelenumi}{\theenumi)}
    \item 
    \item 
    \item 
    \item 
\end{enumerate}

\vspace{0.5cm}
\noindent\textbf{Exercício 11.}
\begin{enumerate}
    \renewcommand{\theenumi}{\alph{enumi}}
    \renewcommand{\labelenumi}{\theenumi)}
    \item 
    \item 
\end{enumerate}

\vspace{0.5cm}
\noindent\textbf{Exercício 12.}
\begin{enumerate}
    \renewcommand{\theenumi}{\alph{enumi}}
    \renewcommand{\labelenumi}{\theenumi)}
    \item 
    \item 
    \item 
    \item 
    \item 
    \item 
    \item 
    \item 
    \item 
    \item 
\end{enumerate}

\vspace{0.5cm}
\noindent\textbf{Exercício 13.}

\vspace{0.5cm}
\noindent\textbf{Exercício 14.}
\begin{enumerate}
    \renewcommand{\theenumi}{\alph{enumi}}
    \renewcommand{\labelenumi}{\theenumi)}
    \item 
    \item 
    \item 
    \item 
\end{enumerate}

\subsection{Função Derivada e Derivadas de Ordem Superior}

\vspace{0.5cm}
\noindent\textbf{Exercício 1.}
\begin{enumerate}
    \renewcommand{\theenumi}{\alph{enumi}}
    \renewcommand{\labelenumi}{\theenumi)}
    \item 
    \item 
    \item 
    \item 
    \item 
    \item 
\end{enumerate}

\vspace{0.5cm}
\noindent\textbf{Exercício 2.}
\begin{enumerate}
    \renewcommand{\theenumi}{\alph{enumi}}
    \renewcommand{\labelenumi}{\theenumi)}
    \item 
    \item 
\end{enumerate}

\vspace{0.5cm}
\noindent\textbf{Exercício 3.}
\begin{enumerate}
    \renewcommand{\theenumi}{\alph{enumi}}
    \renewcommand{\labelenumi}{\theenumi)}
    \item 
    \item 
    \item 
    \item 
\end{enumerate}

\subsection{Notações para a Derivada}

\vspace{0.5cm}
\noindent\textbf{Exercício 1.}
\begin{enumerate}
    \renewcommand{\theenumi}{\alph{enumi}}
    \renewcommand{\labelenumi}{\theenumi)}
    \item 
    \item 
    \item 
    \item 
    \item 
    \item 
    \item 
    \item 
    \item 
    \item 
    \item 
    \item 
    \item 
    \item 
    \item 
    \item 
\end{enumerate}

\vspace{0.5cm}
\noindent\textbf{Exercício 2.}
\begin{enumerate}
    \renewcommand{\theenumi}{\alph{enumi}}
    \renewcommand{\labelenumi}{\theenumi)}
    \item 
    \item 
\end{enumerate}

\vspace{0.5cm}
\noindent\textbf{Exercício 3.}

\vspace{0.5cm}
\noindent\textbf{Exercício 4.}

\vspace{0.5cm}
\noindent\textbf{Exercício 5.}

\vspace{0.5cm}
\noindent\textbf{Exercício 6.}

\vspace{0.5cm}
\noindent\textbf{Exercício 7.}

\vspace{0.5cm}
\noindent\textbf{Exercício 8.}
\begin{enumerate}
    \renewcommand{\theenumi}{\alph{enumi}}
    \renewcommand{\labelenumi}{\theenumi)}
    \item 
    \item 
    \item 
    \item 
    \item 
    \item 
\end{enumerate}

\vspace{0.5cm}
\noindent\textbf{Exercício 9.}

\vspace{0.5cm}
\noindent\textbf{Exercício 10.}

\vspace{0.5cm}
\noindent\textbf{Exercício 11.}

\vspace{0.5cm}
\noindent\textbf{Exercício 12.}

\vspace{0.5cm}
\noindent\textbf{Exercício 13.}

\vspace{0.5cm}
\noindent\textbf{Exercício 14.}

\vspace{0.5cm}
\noindent\textbf{Exercício 15.}

\vspace{0.5cm}
\noindent\textbf{Exercício 16.}
\begin{enumerate}
    \renewcommand{\theenumi}{\alph{enumi}}
    \renewcommand{\labelenumi}{\theenumi)}
    \item 
    \item 
\end{enumerate}

\subsection{Regra da Cadeia para Derivação de Função Composta}
Esta subseção não possui exercícios

\subsection{Aplicações da Regra da Cadeia}

\vspace{0.5cm}
\noindent\textbf{Exercício 1.}
\begin{enumerate}
    \renewcommand{\theenumi}{\alph{enumi}}
    \renewcommand{\labelenumi}{\theenumi)}
    \item 
    \item 
    \item 
    \item 
    \item 
    \item 
    \item 
    \item 
    \item 
    \item 
    \item 
    \item 
    \item 
    \item 
    \item 
    \item 
    \item 
    \item 
    \item 
    \item 
\end{enumerate}

\vspace{0.5cm}
\noindent\textbf{Exercício 2.}

\vspace{0.5cm}
\noindent\textbf{Exercício 3.}

\vspace{0.5cm}
\noindent\textbf{Exercício 4.}
\begin{enumerate}
    \renewcommand{\theenumi}{\alph{enumi}}
    \renewcommand{\labelenumi}{\theenumi)}
    \item 
    \item 
    \item 
    \item 
    \item 
    \item 
    \item 
    \item 
    \item 
    \item 
    \item 
    \item 
    \item 
    \item 
    \item 
    \item 
    \item 
    \item 
    \item 
    \item 
\end{enumerate}

\vspace{0.5cm}
\noindent\textbf{Exercício 5.}
\begin{enumerate}
    \renewcommand{\theenumi}{\alph{enumi}}
    \renewcommand{\labelenumi}{\theenumi)}
    \item 
    \item 
    \item 
    \item 
    \item 
    \item 
    \item 
    \item 
    \item 
    \item 
    \item 
    \item 
    \item 
    \item 
    \item 
    \item 
    \item 
    \item 
    \item 
    \item 
\end{enumerate}

\vspace{0.5cm}
\noindent\textbf{Exercício 6.}

\vspace{0.5cm}
\noindent\textbf{Exercício 7.}

\vspace{0.5cm}
\noindent\textbf{Exercício 8.}

\vspace{0.5cm}
\noindent\textbf{Exercício 9.}

\vspace{0.5cm}
\noindent\textbf{Exercício 10.}

\vspace{0.5cm}
\noindent\textbf{Exercício 11.}

\vspace{0.5cm}
\noindent\textbf{Exercício 12.}

\vspace{0.5cm}
\noindent\textbf{Exercício 13.}

\vspace{0.5cm}
\noindent\textbf{Exercício 14.}

\vspace{0.5cm}
\noindent\textbf{Exercício 15.}
\begin{enumerate}
    \renewcommand{\theenumi}{\alph{enumi}}
    \renewcommand{\labelenumi}{\theenumi)}
    \item 
    \item 
    \item 
    \item 
\end{enumerate}

\vspace{0.5cm}
\noindent\textbf{Exercício 16.}
\begin{enumerate}
    \renewcommand{\theenumi}{\alph{enumi}}
    \renewcommand{\labelenumi}{\theenumi)}
    \item 
    \item 
    \item 
    \item 
    \item 
    \item 
    \item 
    \item 
    \item 
    \item 
    \item 
    \item 
\end{enumerate}

\vspace{0.5cm}
\noindent\textbf{Exercício 17.}

\vspace{0.5cm}
\noindent\textbf{Exercício 18.}

\vspace{0.5cm}
\noindent\textbf{Exercício 19.}

\vspace{0.5cm}
\noindent\textbf{Exercício 20.}

\vspace{0.5cm}
\noindent\textbf{Exercício 21.}

\vspace{0.5cm}
\noindent\textbf{Exercício 22.}

\vspace{0.5cm}
\noindent\textbf{Exercício 23.}
\begin{enumerate}
    \renewcommand{\theenumi}{\alph{enumi}}
    \renewcommand{\labelenumi}{\theenumi)}
    \item 
    \item 
    \item 
\end{enumerate}

\vspace{0.5cm}
\noindent\textbf{Exercício 24.}

\vspace{0.5cm}
\noindent\textbf{Exercício 25.}

\vspace{0.5cm}
\noindent\textbf{Exercício 26.}

\vspace{0.5cm}
\noindent\textbf{Exercício 27.}

\vspace{0.5cm}
\noindent\textbf{Exercício 28.}

\vspace{0.5cm}
\noindent\textbf{Exercício 29.}
\begin{enumerate}
    \renewcommand{\theenumi}{\alph{enumi}}
    \renewcommand{\labelenumi}{\theenumi)}
    \item 
    \item 
\end{enumerate}

\vspace{0.5cm}
\noindent\textbf{Exercício 30.}

\subsection{Derivada de $f(x)^{g(x)}$}

\vspace{0.5cm}
\noindent\textbf{Exercício 1.}
\begin{enumerate}
    \renewcommand{\theenumi}{\alph{enumi}}
    \renewcommand{\labelenumi}{\theenumi)}
    \item 
    \item 
    \item 
    \item 
    \item 
    \item 
    \item 
    \item 
    \item 
    \item 
    \item 
    \item 
    \item 
    \item 
    \item 
    \item 
\end{enumerate}

\vspace{0.5cm}
\noindent\textbf{Exercício 2.}

\vspace{0.5cm}
\noindent\textbf{Exercício 3.}
\begin{enumerate}
    \renewcommand{\theenumi}{\alph{enumi}}
    \renewcommand{\labelenumi}{\theenumi)}
    \item 
    \item 
    \item 
    \item 
    \item 
    \item 
\end{enumerate}

\subsection{Derivação da Função Dada Implicitamente}

\vspace{0.5cm}
\noindent\textbf{Exercício 1.}

\vspace{0.5cm}
\noindent\textbf{Exercício 2.}

\vspace{0.5cm}
\noindent\textbf{Exercício 3.}

\vspace{0.5cm}
\noindent\textbf{Exercício 4.}
\begin{enumerate}
    \renewcommand{\theenumi}{\alph{enumi}}
    \renewcommand{\labelenumi}{\theenumi)}
    \item 
    \item 
    \item 
    \item 
    \item 
    \item 
    \item 
    \item 
    \item 
    \item 
    \item 
    \item 
\end{enumerate}

\vspace{0.5cm}
\noindent\textbf{Exercício 5.}

\vspace{0.5cm}
\noindent\textbf{Exercício 6.}

\vspace{0.5cm}
\noindent\textbf{Exercício 7.}

\vspace{0.5cm}
\noindent\textbf{Exercício 8.}
\begin{enumerate}
    \renewcommand{\theenumi}{\alph{enumi}}
    \renewcommand{\labelenumi}{\theenumi)}
    \item 
    \item 
\end{enumerate}

\vspace{0.5cm}
\noindent\textbf{Exercício 9.}

\vspace{0.5cm}
\noindent\textbf{Exercício 10.}

\vspace{0.5cm}
\noindent\textbf{Exercício 11.}

\subsection{Interpretação de $\frac{dy}{dx}$ como um Quociente Diferencial}

\vspace{0.5cm}
\noindent\textbf{Exercício 1.}
\begin{enumerate}
    \renewcommand{\theenumi}{\alph{enumi}}
    \renewcommand{\labelenumi}{\theenumi)}
    \item 
    \item 
    \item 
    \item 
\end{enumerate}

\vspace{0.5cm}
\noindent\textbf{Exercício 2.}
\begin{enumerate}
    \renewcommand{\theenumi}{\alph{enumi}}
    \renewcommand{\labelenumi}{\theenumi)}
    \item 
    \item 
\end{enumerate}

\vspace{0.5cm}
\noindent\textbf{Exercício 3.}
\begin{enumerate}
    \renewcommand{\theenumi}{\alph{enumi}}
    \renewcommand{\labelenumi}{\theenumi)}
    \item 
    \item 
\end{enumerate}

\vspace{0.5cm}
\noindent\textbf{Exercício 4.}
\begin{enumerate}
    \renewcommand{\theenumi}{\alph{enumi}}
    \renewcommand{\labelenumi}{\theenumi)}
    \item 
    \item 
\end{enumerate}

\subsection{Velocidade e Aceleração. Taxa de Variação}

\vspace{0.5cm}
\noindent\textbf{Exercício 1.}
\begin{enumerate}
    \renewcommand{\theenumi}{\alph{enumi}}
    \renewcommand{\labelenumi}{\theenumi)}
    \item 
    \item 
    \item 
    \item 
\end{enumerate}

\vspace{0.5cm}
\noindent\textbf{Exercício 2.}
\begin{enumerate}
    \renewcommand{\theenumi}{\alph{enumi}}
    \renewcommand{\labelenumi}{\theenumi)}
    \item 
    \item 
    \item 
\end{enumerate}

\vspace{0.5cm}
\noindent\textbf{Exercício 3.}
\begin{enumerate}
    \renewcommand{\theenumi}{\alph{enumi}}
    \renewcommand{\labelenumi}{\theenumi)}
    \item 
    \item 
    \item 
    \item 
\end{enumerate}

\vspace{0.5cm}
\noindent\textbf{Exercício 4.}

\vspace{0.5cm}
\noindent\textbf{Exercício 5.}

\vspace{0.5cm}
\noindent\textbf{Exercício 6.}

\vspace{0.5cm}
\noindent\textbf{Exercício 7.}
\begin{enumerate}
    \renewcommand{\theenumi}{\alph{enumi}}
    \renewcommand{\labelenumi}{\theenumi)}
    \item 
    \item 
    \item 
    \item 
\end{enumerate}

\vspace{0.5cm}
\noindent\textbf{Exercício 8.}
\begin{enumerate}
    \renewcommand{\theenumi}{\alph{enumi}}
    \renewcommand{\labelenumi}{\theenumi)}
    \item 
    \item 
    \item 
    \item 
\end{enumerate}

\vspace{0.5cm}
\noindent\textbf{Exercício 9.}
\begin{enumerate}
    \renewcommand{\theenumi}{\alph{enumi}}
    \renewcommand{\labelenumi}{\theenumi)}
    \item 
    \item 
    \item 
    \item 
    \item 
    \item 
\end{enumerate}

\vspace{0.5cm}
\noindent\textbf{Exercício 10.}
\begin{enumerate}
    \renewcommand{\theenumi}{\alph{enumi}}
    \renewcommand{\labelenumi}{\theenumi)}
    \item 
    \item 
    \item 
    \item 
\end{enumerate}

\vspace{0.5cm}
\noindent\textbf{Exercício 11.}

\vspace{0.5cm}
\noindent\textbf{Exercício 12.}

\vspace{0.5cm}
\noindent\textbf{Exercício 13.}

\vspace{0.5cm}
\noindent\textbf{Exercício 14.}
\begin{enumerate}
    \renewcommand{\theenumi}{\alph{enumi}}
    \renewcommand{\labelenumi}{\theenumi)}
    \item 
    \item 
\end{enumerate}

\vspace{0.5cm}
\noindent\textbf{Exercício 15.}

\vspace{0.5cm}
\noindent\textbf{Exercício 16.}

\vspace{0.5cm}
\noindent\textbf{Exercício 17.}

\vspace{0.5cm}
\noindent\textbf{Exercício 18.}

\vspace{0.5cm}
\noindent\textbf{Exercício 19.}

\vspace{0.5cm}
\noindent\textbf{Exercício 20.}

\vspace{0.5cm}
\noindent\textbf{Exercício 21.}

\subsection{Problemas Reta Tangente e Normal ao Gráfico de uma Função}

\vspace{0.5cm}
\noindent\textbf{Exercício 1.}
\begin{enumerate}
    \renewcommand{\theenumi}{\alph{enumi}}
    \renewcommand{\labelenumi}{\theenumi)}
    \item 
    \item 
    \item 
    \item 
\end{enumerate}

\vspace{0.5cm}
\noindent\textbf{Exercício 2.}

\vspace{0.5cm}
\noindent\textbf{Exercício 3.}

\vspace{0.5cm}
\noindent\textbf{Exercício 4.}

\vspace{0.5cm}
\noindent\textbf{Exercício 5.}

\vspace{0.5cm}
\noindent\textbf{Exercício 6.}
\begin{enumerate}
    \renewcommand{\theenumi}{\alph{enumi}}
    \renewcommand{\labelenumi}{\theenumi)}
    \item 
    \item 
\end{enumerate}

\vspace{0.5cm}
\noindent\textbf{Exercício 7.}

\vspace{0.5cm}
\noindent\textbf{Exercício 8.}

\vspace{0.5cm}
\noindent\textbf{Exercício 9.}

\vspace{0.5cm}
\noindent\textbf{Exercício 10.}

\vspace{0.5cm}
\noindent\textbf{Exercício 11.}

\vspace{0.5cm}
\noindent\textbf{Exercício 12.}

\vspace{0.5cm}
\noindent\textbf{Exercício 13.}

\vspace{0.5cm}
\noindent\textbf{Exercício 14.}

\vspace{0.5cm}
\noindent\textbf{Exercício 15.}

\subsection{Exercícios do Capítulo}

\vspace{0.5cm}
\noindent\textbf{Exercício 1.}
\begin{enumerate}
    \renewcommand{\theenumi}{\alph{enumi}}
    \renewcommand{\labelenumi}{\theenumi)}
    \item 
    \item 
    \item 
    \item 
    \item 
    \item 
    \item 
    \item 
    \item 
\end{enumerate}

\vspace{0.5cm}
\noindent\textbf{Exercício 2.}
\begin{enumerate}
    \renewcommand{\theenumi}{\alph{enumi}}
    \renewcommand{\labelenumi}{\theenumi)}
    \item 
    \item 
    \item 
    \item 
    \item 
    \item 
    \item 
    \item 
    \item 
    \item 
    \item 
    \item 
    \item 
    \item 
    \item 
    \item 
    \item 
    \item 
    \item 
    \item 
\end{enumerate}

\vspace{0.5cm}
\noindent\textbf{Exercício 3.}
\begin{enumerate}
    \renewcommand{\theenumi}{\alph{enumi}}
    \renewcommand{\labelenumi}{\theenumi)}
    \item 
    \item 
    \item 
    \item 
\end{enumerate}

\vspace{0.5cm}
\noindent\textbf{Exercício 4.}

\vspace{0.5cm}
\noindent\textbf{Exercício 5.}

\vspace{0.5cm}
\noindent\textbf{Exercício 6.}

\vspace{0.5cm}
\noindent\textbf{Exercício 7.}

\vspace{0.5cm}
\noindent\textbf{Exercício 8.}

\vspace{0.5cm}
\noindent\textbf{Exercício 9.}

\vspace{0.5cm}
\noindent\textbf{Exercício 10.}

\vspace{0.5cm}
\noindent\textbf{Exercício 11.}

\vspace{0.5cm}
\noindent\textbf{Exercício 12.}

\vspace{0.5cm}
\noindent\textbf{Exercício 13.}

\vspace{0.5cm}
\noindent\textbf{Exercício 14.}

\vspace{0.5cm}
\noindent\textbf{Exercício 15.}

\vspace{0.5cm}
\noindent\textbf{Exercício 16.}

\vspace{0.5cm}
\noindent\textbf{Exercício 17.}

\vspace{0.5cm}
\noindent\textbf{Exercício 18.}

\vspace{0.5cm}
\noindent\textbf{Exercício 19.}

\vspace{0.5cm}
\noindent\textbf{Exercício 20.}
\begin{enumerate}
    \renewcommand{\theenumi}{\alph{enumi}}
    \renewcommand{\labelenumi}{\theenumi)}
    \item 
    \item 
\end{enumerate}

\vspace{0.5cm}
\noindent\textbf{Exercício 21.}
\begin{enumerate}
    \renewcommand{\theenumi}{\alph{enumi}}
    \renewcommand{\labelenumi}{\theenumi)}
    \item 
    \item 
\end{enumerate}

\vspace{0.5cm}
\noindent\textbf{Exercício 22.}
\begin{enumerate}
    \renewcommand{\theenumi}{\alph{enumi}}
    \renewcommand{\labelenumi}{\theenumi)}
    \item 
    \item 
    \item 
    \item 
\end{enumerate}

\vspace{0.5cm}
\noindent\textbf{Exercício 23.}

\vspace{0.5cm}
\noindent\textbf{Exercício 24.}

\vspace{0.5cm}
\noindent\textbf{Exercício 25.}
\begin{enumerate}
    \renewcommand{\theenumi}{\alph{enumi}}
    \renewcommand{\labelenumi}{\theenumi)}
    \item 
    \item 
\end{enumerate}

\vspace{0.5cm}
\noindent\textbf{Exercício 26.}
\begin{enumerate}
    \renewcommand{\theenumi}{\alph{enumi}}
    \renewcommand{\labelenumi}{\theenumi)}
    \item 
    \item 
\end{enumerate}

\vspace{0.5cm}
\noindent\textbf{Exercício 27.}
\begin{enumerate}
    \renewcommand{\theenumi}{\alph{enumi}}
    \renewcommand{\labelenumi}{\theenumi)}
    \item 
    \item 
\end{enumerate}

\vspace{0.5cm}
\noindent\textbf{Exercício 28.}
\begin{enumerate}
    \renewcommand{\theenumi}{\alph{enumi}}
    \renewcommand{\labelenumi}{\theenumi)}
    \item 
    \item 
\end{enumerate}

\vspace{0.5cm}
\noindent\textbf{Exercício 29.}
\begin{enumerate}
    \renewcommand{\theenumi}{\alph{enumi}}
    \renewcommand{\labelenumi}{\theenumi)}
    \item 
    \item 
\end{enumerate}

\vspace{0.5cm}
\noindent\textbf{Exercício 30.}

\vspace{0.5cm}
\noindent\textbf{Exercício 31.}

\vspace{0.5cm}
\noindent\textbf{Exercício 32.}

\vspace{0.5cm}
\noindent\textbf{Exercício 33.}
\begin{enumerate}
    \renewcommand{\theenumi}{\alph{enumi}}
    \renewcommand{\labelenumi}{\theenumi)}
    \item 
    \item 
\end{enumerate}

\vspace{0.5cm}
\noindent\textbf{Exercício 34.}

\vspace{0.5cm}
\noindent\textbf{Exercício 35.}

\vspace{0.5cm}
\noindent\textbf{Exercício 36.}
\begin{enumerate}
    \renewcommand{\theenumi}{\alph{enumi}}
    \renewcommand{\labelenumi}{\theenumi)}
    \item 
    \item 
\end{enumerate}

\vspace{0.5cm}
\noindent\textbf{Exercício 37.}

\vspace{0.5cm}
\noindent\textbf{Exercício 38.}

\vspace{0.5cm}
\noindent\textbf{Exercício 39.}
\begin{enumerate}
    \renewcommand{\theenumi}{\alph{enumi}}
    \renewcommand{\labelenumi}{\theenumi)}
    \item 
    \item 
    \item 
    \item 
\end{enumerate}

\vspace{0.5cm}
\noindent\textbf{Exercício 40.}

\vspace{0.5cm}
\noindent\textbf{Exercício 41.}
\begin{enumerate}
    \renewcommand{\theenumi}{\alph{enumi}}
    \renewcommand{\labelenumi}{\theenumi)}
    \item 
    \item 
    \item 
    \item 
    \item 
    \item 
\end{enumerate}

\vspace{0.5cm}
\noindent\textbf{Exercício 42.}

\vspace{0.5cm}
\noindent\textbf{Exercício 43.}

\vspace{0.5cm}
\noindent\textbf{Exercício 44.}

\vspace{0.5cm}
\noindent\textbf{Exercício 45.}

\vspace{0.5cm}
\noindent\textbf{Exercício 46.}

\vspace{0.5cm}
\noindent\textbf{Exercício 47.}

\vspace{0.5cm}
\noindent\textbf{Exercício 48.}

\vspace{0.5cm}
\noindent\textbf{Exercício 49.}

\end{document}
